\documentclass[11pt]{article}
    \usepackage[utf8]{inputenc}
    \usepackage[T1]{fontenc}
    \usepackage[margin=1in]{geometry}
    \usepackage{hyperref}
    \usepackage{enumitem}
    \title{3D Design in Tinkercad, Introductory Lesson. Design your own Minecraft Bobblehead}
    \author{Piter Garcia}
    \date{March 20, 2026}
    \begin{document}
    \maketitle
    \section*{Quick Scan}
    \begin{itemize}[leftmargin=*]
    \item Grade: Grade 5
    \item Workbook sessions: Session 21
    \item Class blocks: Day 2 10:45-11:35 Pickett, Day 3 10:45-11:35 McGlashon, Day 5 10:45-11:35 Pecora
    \item Reference file: 3d-design-in-tinkercad-introductory-lesson-design-your-own-minecraft-bob.bib
    \end{itemize}
    \section*{School Planning Goals and Inclusion}
    \begin{itemize}[leftmargin=*]
    \item What is expected: I can describe the steps taken and choices made to design and develop a Bobblehead (4-6.CT.10) | I can decompose a problem into smaller named tasks (4-6.CT.4)
    \item What we achieved: This lesson points to local transcript or evidence files in the workspace so the packet can be revised from what actually happened in class.
    \item What needs improvement: stronger proof, cleaner assessment notes, and tighter lesson artifacts.
    \item How to improve: add transcript proof, one saved artifact, and clearer success criteria.
    \item Inclusion focus: use visual modeling, chunked directions, predictable routines, and flexible participation roles.
    \end{itemize}
    \section*{Official References}
    \begin{itemize}[leftmargin=*]
    \item \url{https://csteachers.org/k12standards/interactive/}
\item \url{https://education.minecraft.net/en-us/resources/computer-science/coding-with-minecraft}
\item \url{https://www.tinkercad.com/}

    \end{itemize}
    \end{document}
